\documentclass[11pt]{article}

\usepackage[margin=1.05in]{geometry}
\usepackage{amsmath,amssymb,amsthm}
\usepackage{booktabs}
\usepackage{microtype}
\usepackage{xcolor}
\usepackage[colorlinks=true,linkcolor=blue!60!black,citecolor=blue!60!black,urlcolor=blue!60!black]{hyperref}

\newcommand{\PT}{\mathcal{PT}}
\newcommand{\ee}{\mathrm{e}}
\newcommand{\ii}{\mathrm{i}}

\newtheorem{proposition}{Proposition}
\newtheorem{claim}{Claim}
\newtheorem{conjecture}{Conjecture}
\theoremstyle{remark}
\newtheorem{remark}{Remark}

\title{Ghosts as an interpretation experiment:\\
what indivisible processes say to Everett}
\author{Jonathan Torkelson\\[2pt]
\small Embeddetech, Inc., Tulsa, Oklahoma, USA \quad
\texttt{jtorkelson@embeddetech.com}}
\date{July 19, 2026}

\begin{document}
\maketitle

\begin{abstract}
\noindent\textbf{Draft --- not for circulation.}
The quantum interpretation debate is usually left to philosophy because
the leading readings are believed empirically equivalent. This paper
argues that the equivalence is conditional on the matter content of the
world, and makes three claims of decreasing modesty about what the
indivisible-stochastic-process reconstruction of quantum theory, and
the ghost sectors studied in its companion papers, say to the Everett
interpretation. \emph{First}, Barandes' reconstruction (the wave
function as nomic, law-like bookkeeping for an indivisible stochastic
process on configurations) removes the \emph{no-alternative} argument
for $\psi$-ontology without refuting it, and the Bender--Mannheim
program (the
Hilbert-space inner product derived from the operator algebra rather
than postulated) sharpens the resulting underdetermination by
dissolving a second postulated primitive; the Pusey--Barrett--Rudolph
theorem does not apply, by the theorem's own scope---an exit we take
from Barandes and cite rather than establish, so this first claim is
conditional on its soundness; the exit's real load is the nomic
reclassification of $\psi$. \emph{Second}, the
ghost trilemma of the companion papers---coupled ghosts, empirical
Lorentz invariance, definite configurations: any two, never all three,
at tree level---\emph{locates} Everett as the
drop-definite-configurations horn; this is a location, not a
criticism, since manifest covariance is exactly what dropping
configurations buys. \emph{Third}---the paper's own claim, and it is
conjectural---a coupled ghost sector would separate the two ontological classes
\emph{observationally with the Born rule intact}: a beables ontology
is foliation-anchored, the foliation's determinacy is a measurable
property of the state that is generically nonzero on
ordinary matter
but whose observable manifestation is shielded there, and a ghost
sector would (conditionally) unshield it through a frame-locked
kinematic dipole and spectral edge, while no-preferred-slicing
ontologies do not entail such a signal, remaining compatible with either
outcome. The claim rides on the converse of a
no-signaling shield proven only for a finite core, and it
discriminates a class (foliation-anchored single-world versus
no-slicing), not a single interpretation. It separates the classes
cleanly only when the ghost Lagrangian is held fixed and
Lorentz-invariant and the individuating frame arises solely as a
state-functional foliation: a fundamentally Lorentz-violating
completion, or a birth-surface (cosmological initial-condition)
structure, is available to a no-slicing interpreter too and would
reproduce the same signal---the birth surface at no cost, the
Lorentz-violating completion at the price of manifest covariance---so
a positive detection is
merely \emph{consistent with} the foliation-anchored class---being
confounded by such completions or initial conditions---while the null
direction is the cleaner discriminator,
because the beables class \emph{entails} the signal (of the computed
template amplitude, for an established sector, under the assumed
comoving-frame identification---a realization off that frame, below
the template's signal fraction, or individuated by a completion other
than the template's frame cutoff, survives the pointed null) whereas a
no-slicing ontology merely \emph{permits} it. In fact no purely state-functional rescue of
a ghost vacuum has yet been exhibited, so on current knowledge every
concrete realization of the positive signal lies in the confounded
branch and the clean case is conjectural over and above the
shield-converse conjecture; this only strengthens the null direction,
since then every viable beables completion predicts a signal---at the
template's sensitivity, for a template-saturating realization. The
closest ancestor of the escape,
Valentini's quantum non-equilibrium, requires Born-rule violations;
here the discriminator is a sector, not a rule failure. (One other
equilibrium route to foliation observability, via contested
arrival-time observables, is disjoint from ours.) Along the way
we state the bookkeeping objection an Everettian should raise against
the stochastic reconstruction---to our knowledge for the first time in
print---and give a partly conditional answer.
\end{abstract}

\newpage
\begin{quote}
\small\noindent\textbf{Note on provenance.} This manuscript is an
experiment in AI-conducted research as much as a physics paper. The
project began when the author asked the large language model Claude
Fable~5 (Anthropic): \emph{``Is it not possible that the antilinear
PT/time-reversal symmetry element could maybe predict indivisible
non-Markovian processes?''} Everything that followed was performed by
the model: the research program, the theorems and proofs, the numerical
experiments, the literature search, and this text. The model was
repeatedly prompted to attack its own results adversarially and to
search for prior work. The author, a professional engineer rather than
a physicist, directed and reviewed the sessions but makes no claim to
comprehension of the findings, and is circulating the manuscript for
independent review by professional physicists. Accordingly, nothing
here should be taken on trust: every claim is argued in full, and the
originating conversation, project briefing, and full research notebook
are preserved in the program's repositories \cite{indivisible}. This
paper is a philosophy-of-physics synthesis: it proves no new theorem,
and its central claim is explicitly conditional.
\end{quote}

%=====================================================================
\section{Introduction}
\label{sec:intro}

The interpretation debate in quantum mechanics is usually adjudicated
by philosophy, because the contenders are believed to be empirically
equivalent: Everettians, Bohmians, and their relatives agree on every
Born-rule expectation value, so the argument proceeds by theoretical
virtue---parsimony, covariance, explanatory depth---rather than by
measurement. Where the data are silent, taste decides. This paper
argues that the silence is not unconditional. It is conditional on the
matter content of the world, and there is a specific, physically
motivated sector---a ghost sector, in the sense of higher-derivative
and negative-norm quantum field theory---whose presence would break
the equivalence between two broad classes of interpretation
\emph{while leaving the Born rule intact everywhere}. Interpretations
that anchor definite configurations to a spacelike slicing would
predict a frame-locked observational signature; interpretations with
no preferred slicing would not entail it, remaining compatible with
either outcome. The experiment may
never become available: nature may simply contain no coupled ghost
sector, and one of the two supporting results below is a conjecture
rather than a theorem. But the conditional itself, we will argue, is
enough to change what the standing philosophical arguments can
establish.

The paper makes three claims, in decreasing order of modesty---each
weaker than a reader might fear.

\begin{enumerate}
\item \textbf{Underdetermine.} Barandes' stochastic--quantum
correspondence, in which the wave function is nomic, law-like
bookkeeping for an indivisible stochastic process on configurations
\cite{BarandesPoP,BarandesTheorem}, removes the \emph{no-alternative}
argument for $\psi$-ontology; the Bender--Mannheim antilinearity
program, in which the Hilbert-space inner product is derived from the
operator algebra rather than postulated
\cite{BenderMannheim2008,Mannheim2020}, reinforces this by dissolving a
second postulated primitive and thereby sharpening the
underdetermination. They do not refute the Everett interpretation; they
restore an underdetermination that the no-alternative argument claimed
to have closed. This claim is
conditional on the soundness of Barandes' exit from the PBR theorem,
which we cite but do not here establish
(Sections~\ref{sec:landscape} and~\ref{sec:underdet}).
\item \textbf{Locate.} The ghost trilemma established in the
program's companion papers \cite{instantaneous,choosing}---coupled
ghosts, empirical Lorentz invariance, definite configurations: any
two, never all three, at tree level---\emph{locates} Everett as the
horn that drops definite configurations. This is a placement on a
map, not an objection: manifest covariance, the Everettian's prized
possession, is exactly what dropping configurations buys
(Section~\ref{sec:trilemma}).
\item \textbf{Separate.} A coupled ghost sector would separate
foliation-anchored single-world ontologies from no-preferred-slicing
ontologies \emph{observationally, with the Born rule intact}: the
beable kinematics of such a sector is frame-locked, and the frame
would be visible as a kinematic dipole and a direction-modulated
spectral edge in the sector's decay products
\cite{instantaneous}. This third claim is \emph{conjectural}: it is
the converse of a no-signaling shield that is proven only for a
finite-configuration core and conjectured at field-theory strength,
and we state it conditionally throughout
(Section~\ref{sec:separation}). It carries a further restriction: the
separation is clean only when the ghost Lagrangian is held fixed and
\emph{Lorentz-invariant} and the individuating frame arises as a
state-functional foliation. If the frame instead enters through a
fundamentally Lorentz-violating completion or a birth-surface
(cosmological initial-condition) structure, a no-slicing interpreter
can adopt the same structure and reproduce the same signal---at no cost
in the birth-surface case, at the price of manifest covariance in the
Lorentz-violating one---so the positive direction is confounded and the
null direction is the
cleaner discriminator (Section~\ref{sec:separation}). No purely state-functional
rescue has in fact been exhibited, so on current knowledge the clean
case is conjectural over and above the shield-converse conjecture, and
every concrete realization of the positive signal lies in the
confounded branch---which only strengthens the null direction
(Section~\ref{sec:separation}).
\end{enumerate}

One paragraph of orientation about tone. The Everett interpretation's
central virtue is real and this paper trades on it rather than
disputing it: Everett takes the formalism literally, adds nothing to
it, and in return gets manifest relativistic covariance and a clean
story about why no experiment has ever needed a preferred frame
\cite{Wallace2012}. Nothing below argues that Everett is wrong.
The paper's business is pricing, not sneering: every horn of the
trilemma has a price, Everett's included, and our contribution is to
say what the prices are and---the new part---to identify the register
in which nature could be induced to express an opinion.

A structural guard, stated up front because the argument depends on
it. Barandes asserts in print that his formulation has \emph{no
fundamentally preferred foliation} \cite{BarandesPoP} and is causally
local \cite{BarandesCausal,BHK2025}. The foliation-anchoring on which our
third claim rests must therefore be---and is---a \emph{conditional
extension} of his framework, not a contradiction of it. His claims,
on the reading we defend in Section~\ref{sec:separation}, concern the
causal and signaling structure of ordinary matter sectors,
where we agree the foliation is hidden; our claim is that a
\emph{ghost} sector would promote the beable-kinematics foliation from
convention to observable. We extend his framework into a regime it
does not treat; we do not disagree with it about the regime it does.
On ordinary matter the two positions are indistinguishable, which is
precisely the point (Section~\ref{sec:separation}).

Section~\ref{sec:landscape} reviews the two reconstructions and the
Pusey--Barrett--Rudolph question. Section~\ref{sec:underdet} states
what the resulting underdetermination shows and what it does not, and
formulates and answers the bookkeeping objection.
Section~\ref{sec:trilemma} restates the trilemma and places Everett on
it. Section~\ref{sec:separation} develops the separation claim and its
conditional status. Section~\ref{sec:limits} lists what does not
follow, and Section~\ref{sec:discussion} concludes.

%=====================================================================
\section{The reconstruction landscape}
\label{sec:landscape}

\subsection{Barandes: $\psi$ as nomic bookkeeping}
\label{sec:barandes}

The stochastic--quantum correspondence
\cite{BarandesPoP,BarandesTheorem} represents any unitary quantum
system as a stochastic process on a configuration space $\mathcal{C}$:
the system determinately occupies one configuration at each time, and
the theory supplies time-dependent transition probabilities
$\Gamma_{ij}(t)$, obtained from the unitary $U(t)$ as
$\Gamma_{ij}(t)=|U_{ij}(t)|^{2}$ in the configuration basis. The
resulting matrices are \emph{unistochastic}, and---this is the
load-bearing structural feature---generically \emph{indivisible}: the
process specifies transition probabilities from the initial time to
each later time, but these do not factor through intermediate times as
a Markov chain would,
$\Gamma(t)\neq\Gamma(t\mid t')\,\Gamma(t')$ for stochastic
intermediaries, except at isolated \emph{division events}.
Indivisibility replaces Markovianity; the process has memory, but the
memory is carried by the law, not stored in any additional beable.

On this reading the wave function is demoted from object to
instrument. Configurations are the ontology; $\psi$ is the compendium
of conditional information the law needs in order to issue transition
probabilities---in Barandes' published formulation, the wave function
``has a law-like or nomic character'' \cite{BarandesPoP}. We note a
terminological drift in the primary sources that we track but do not
paper over: the published account of 2025 speaks of the nomic wave
function \cite{BarandesPoP}, while the 2026 development recasts the
same structural point in latent-variable and hidden-Markov-model
language \cite{BarandesHMM}. The gloss we will use in this paper---%
$\psi$ as \emph{derived bookkeeping}, the memory record of an
indivisible process---is our own, and we flag it as such; nothing
below depends on more than the shared structural content of the two
framings, namely that $\psi$ is not a beable.

\paragraph{PBR does not apply---by the theorem's own scope.}
The standard reflex against any $\psi$-nonontic reading is the
Pusey--Barrett--Rudolph theorem \cite{PBR}. But PBR is a theorem
\emph{inside} the ontological-models framework of Harrigan and
Spekkens \cite{HS2010}, and that framework builds in a presupposition
the indivisible reconstruction denies. Harrigan and Spekkens assume
that every system has an underlying \emph{ontic state} $\lambda$ and
that each preparation fixes a probability distribution over $\lambda$;
PBR is then a single-time statement about such distributions, with
measurement responses depending on the $\lambda$ sampled at
preparation. (We are careful here: Harrigan and Spekkens draw no
nomic or law-like carve-out of their own---their classification
situates every $\psi$ within the ontic-state framework (as ontic, or as
epistemic about $\lambda$), and they classify
de Broglie--Bohm's $\psi$ as ontic. Their Section~6 does note, following
a suggestion of Wiseman, that a wave function ``deemed to be nomic, that
is, law-like'' would sit outside the ontic/epistemic dichotomy---but they
raise it only as a possibility ``deserving of further scrutiny,'' not as
a licensed exit they endorse.) An indivisible, non-Markovian process is not \emph{naturally} of the
framework's form: its response at a later time is not fixed by any
single-time ontic state carried forward from preparation, because the
dynamics has memory that is not stored in a time-local $\lambda$. We
say ``memory'' rather than repeating the earlier gloss of a
``Markovian assumption about the ontology'': PBR does not itself
assume a Markov chain, but its ontic states are single-time objects,
and memory across time is what an indivisible process has and such a
$\lambda$ does not encode. This last step invites a rejoinder we should
name: Harrigan and Spekkens place no formal restriction on what
$\lambda$ may be, so a modeler could take $\lambda$ to be a whole
trajectory segment, or to include $\psi$ outright, whereupon PBR applies
and returns the onticity of \emph{that} enriched $\lambda$. Against such
embeddings the exit does not rest on single-time-ness at all; its real
load is the \emph{nomic reclassification}---that $\psi$ belongs to the
law rather than to the ontic state---which is a classification claim,
not a single-time no-go. Placing indivisible dynamics outside the
framework on these grounds is Barandes' claim, not ours; he states it
explicitly in the published version of the correspondence paper---%
indivisible quantum theory is based on non-Markovian stochastic
dynamics and so lies outside the ontological-models framework, the
wave function being ``neither ontic nor entirely epistemic'' but
nomic, and consequently the PBR theorem ``does not apply''
\cite{BarandesPoP}. We cite and rely on this exit; we do not
independently establish it, and claim (i) is conditional on its
soundness. We have not found a published challenge to the adequacy of
the exit; stress-testing it is work this paper flags as open but does
not undertake.

Two remarks keep the exit honest. First, it is not a loophole of
convenience---though this establishes less than it may seem. The
non-Markovianity is not a modeling choice adopted to dodge a theorem:
it is a structural, computed property of the dynamics. The program's
foundational companion \cite{indivisible} develops the operative
test---negativity of the unique divisor when the earlier transition
matrix is invertible---and proves closed-form results about it,
including a terminal-indivisibility theorem with an exact generator
floor in the broken-$\PT$ phase. Indivisibility is a
convention-independent sign condition one calculates, not a knob one
turns. But we must keep two claims apart. Structural-ness shows only
that the exit is \emph{non-arbitrary}---the non-Markovianity is a
genuine feature of the dynamics, not an ad hoc device introduced to
evade PBR. It does \emph{not} by itself establish that non-Markovianity
places the theory outside PBR's scope; that further step is exactly the
(cited, not here established) scope claim of the previous paragraph,
and it should not be smuggled in on the back of the non-arbitrariness
point. Second, the exit runs through the nomic, non-Markovian route,
\emph{not} through denying preparation independence. Gao's variant,
which recovers single-system $\psi$-onticity after removing the
preparation-independence assumption \cite{PBRnoPI}, therefore tightens
that framework's internal assumptions without moving the framework's
boundary; it is important for the taxonomy of $\psi$-epistemic models
inside the framework, and irrelevant to a reading that was never inside
it. We make this point about Gao's preparation-independence-removing
route specifically, not about PBR strengthenings in general.

\subsection{Mannheim: the inner product as output}
\label{sec:mannheim}

The second reconstruction attacks a different primitive. In standard
quantum mechanics the Hilbert-space inner product---the object that
turns amplitudes into probabilities---is postulated. In the
$\PT$-symmetric and antilinear quantization program of Bender and
Mannheim, Hermiticity of the Hamiltonian is replaced by invariance
under an antilinear symmetry, and the inner product is
\emph{constructed} from the dynamics: one solves for the metric
operator that renders the evolution unitary, and the probability
interpretation is an output of the operator algebra rather than an
input to it. The flagship result is the no-ghost theorem for the
fourth-order Pais--Uhlenbeck oscillator \cite{BenderMannheim2008}:
a theory whose naive inner product exhibits negative norms is, in the
correct dynamically determined inner product, a unitary theory with a
positive spectrum. The program extends from quantum mechanics through
Pauli--Villars regularization to fourth-order quantum gravity
\cite{Mannheim2020}.

For the interpretation debate the lesson is orthogonal to ghost
physics: the inner product is not sacred. Combined with the
stochastic--quantum correspondence, the two pieces of the quantum
formalism most often taken as ontological primitives---the state
vector and the probability-bearing inner product---are \emph{each}
reconstructible from something else: the state from a stochastic
process on configurations, the inner product from the demand of
probability conservation on the operator algebra. The program's
foundational paper \cite{indivisible} shows further that the two
reconstructions are compatible and interlocking: unbroken antilinear
symmetry is exactly the condition for a probability-conserving
process to exist without enlarging the configuration space, and the
metric-corrected process is the ordinary (generically indivisible)
process of the isospectral Hermitian partner. The underdetermination
argued in the next section is therefore sharpened twice over: not one
alternative reading but a pair of independent reconstructions, each
dissolving a different primitive.

%=====================================================================
\section{What underdetermination shows, and what it does not}
\label{sec:underdet}

The Everett interpretation's strongest argument has always been an
inference to the only explanation. In the form given by Deutsch
\cite{Deutsch1996}, sharpened by Brown and Wallace
\cite{BrownWallace}, and developed systematically by Wallace
\cite{Wallace2012}, it runs: take the bare quantum formalism
literally---a quantum state evolving unitarily---and you get a
multiverse; every alternative either adds structure the formalism
does not need (corpuscles, a pilot wave) or mutilates the dynamics
(collapse); therefore realism about quantum mechanics \emph{is}
Everettianism. The argument's force lives entirely in its premise:
that nothing else works without ad hoc additions.

The reconstructions of Section~\ref{sec:landscape} break the premise,
not the conclusion. An ontology of definite configurations governed by
a law issuing indivisible transition probabilities reproduces every
Born-rule prediction; no \emph{beable} is added to the formalism---%
rather, two things usually postulated are derived. This is not quite
``nothing is added,'' and we should not claim it is: the stochastic
reading does carry a commitment the bare formalism leaves open, namely
a preferred configuration basis, since Barandes' process is defined in
the configuration basis. We take this commitment to be benign---a
preferred basis has to be selected on any reading that recovers
definite outcomes, and Everett's own recovery leans on a
decoherence-selected basis whose emergence is not obviously cheaper,
resting as it does on a system--environment split it does not
derive---but we flag it rather than bury it. The wave function is not an extra beable riding
alongside the configurations (the disanalogy with pilot-wave theory is
developed below); it is the law's bookkeeping. And the inner product
that makes the probabilities is itself an output. If the
reconstructions are coherent---and they have both defenders and
critics, the no-ghost/benign-ghost resolution of the Pais--Uhlenbeck
sector in particular being contested in the published literature, so we
take only the weaker premise that they are coherent \emph{enough} to
constitute a live alternative---then ``take the formalism literally''
no longer picks out a unique ontology, because the formalism admits
more than one literal reading. Everett remains exactly as consistent as
before; what it loses is the claim to be the \emph{only} non-mutilating
realism. Here a distinction matters: Deutsch's version of the argument
is \emph{absolute}---there is no alternative---while Wallace's is
\emph{comparative}, that Everett is the best-costed realism
\cite{Deutsch1996,Wallace2012}. The reconstructions defeat the
absolute premise outright; against the comparative version they reopen
a cost comparison rather than win it. Underdetermination is restored,
not resolved.

\paragraph{The bookkeeping objection, stated.}
There is an obvious and, we think, obligatory Everettian reply, which
to our knowledge has not been put in print against the stochastic
reconstruction; we state it in its strongest form because our answer
to it organizes the rest of the paper. A word up front on the
structure of that answer, since it bears on how much weight the answer
can carry. Of the three replies below, the first two---parity and the
disanalogy with pilot-wave theory---are defensive and available
unconditionally; the third, which is the only one that would
positively convert the dispute into a factual disagreement, is the
paper's conjectural claim (iii), conditional on both a coupled ghost
sector existing and the shield-converse conjecture holding. In the
ghost-free regime, therefore, the objection is met only by parity and
the disanalogy---contestable, defensive replies---and is not decisively
refuted. The objection runs: \emph{the
indivisible process is bookkeeping on the universal wave function, not
an alternative to it}. The transition matrices are manufactured from
the unitary, $\Gamma_{ij}=|U_{ij}|^{2}$; the process's celebrated
memory is $\psi$'s memory; the division events are $\psi$'s
interference nodes. Calling the configurations ``beables'' renames the
multiverse rather than removing it---one has merely elected to track a
single thread through a structure whose full contents the law still
carries. This is Deutsch's charge against pilot-wave theory---%
``parallel-universes theories in a state of chronic denial''
\cite{Deutsch1996}---transposed to the stochastic key, with the
Brown--Wallace development \cite{BrownWallace} supplying the general
form: if the branching structure is present in the theory's dynamical
core, dressing one branch with ontological significance is an
addition of words, not of physics.

\paragraph{The answer.} Three points, in increasing order of weight.

\emph{(i) Parity.} Bookkeeping arguments cut both ways between
mathematically interconvertible formulations. The Everettian says the
process is bookkeeping on $\psi$; the stochastic theorist says
$\psi$ is bookkeeping on the process---and the correspondence
theorems license both directions of translation
\cite{BarandesTheorem}. When two formulations are interconvertible,
the assertion ``yours is merely a re-description of mine'' is
symmetric under exchange of the disputants, and therefore cannot by
itself settle whose structure is ontological. It is worth adding that
even the empirical-equivalence proofs underlying such translation
arguments are less watertight than folklore holds: Manero, Muci\~no,
and Okon have recently shown that the standard equivalence claims for
pilot-wave theory rest on a disputable auxiliary assumption
\cite{MMO2025}. Equivalence is a theorem about specified regimes, not
a metaphysical axiom.

\emph{(ii) The disanalogy with pilot-wave theory.} Deutsch's charge
had bite against Bohmian mechanics for a specific structural reason:
that theory posits an \emph{ontic} wave function \emph{and}
corpuscles. The ontic $\psi$-field contains the full branching
structure as physical content; the corpuscles merely index one
branch; the ``empty'' branches persist as real, structured, inhabited
by functionalism's lights, and the denial charge is the observation
that they are still there. Valentini's rebuttal \cite{Valentini0811}
and Brown's counter-rebuttal \cite{Brown0901} contest whether the
empty branches deserve the name ``worlds''; the debate is live
precisely because the pilot wave is ontic. The stochastic
reconstruction has no analogous structure to be in denial about.
There is no ontic $\psi$ riding alongside the configurations: the
wave function is nomic---an aspect of the law---and a law is not a
container of uninhabited worlds any more than a Hamiltonian is. To
insist that the branching structure ``is still there because it is in
the law'' is to assume that whatever appears in the law's bookkeeping
is thereby physically real---which is exactly the $\psi$-ontology in
dispute, now premised rather than argued. The objection, pressed to
this point, threatens to beg the question. We concede, however, that
this disanalogy is contestable rather than decisive. A Wallace-style
real-patterns functionalist \cite{Wallace2012} can resist the move
from ``nomic'' to ``no worlds'': if branch structure is a real
pattern, it may supervene on whatever carries that pattern---including
a law---so that locating the branching in $\psi$'s nomic role need not
by itself drain it of worlds. Our reply is not that this reading is
incoherent but that it must \emph{argue} for law-hosted inhabited
worlds rather than read them off the analogy with an ontic field: ``a
law is not a container of worlds'' is a prima facie disanalogy that
shifts the burden, not a proof that discharges it. We rest the
defensive case on that shift, not on closing the question.

\emph{(iii) The point that could be decisive is empirical---but only
conditionally.} Parity and disanalogy
are defensive; the third point would not be, \emph{if} it obtains. If
two readings of a formalism can be brought to disagree about the
expectation value of an observable in \emph{some} physically
specifiable sector, then neither is bookkeeping for the other---%
bookkeeping cannot contradict what it keeps books on. They would then
be different theories that coincide on the sectors so far observed.
Section~\ref{sec:separation} argues, conditionally, that the
foliation-anchored and no-slicing readings are of this kind: a coupled
ghost sector would elicit different predictions from them with the Born
rule intact on both sides. But this reply is only as strong as its
conditions: it requires (1)~a coupled ghost sector to exist, (2)~%
Conjecture~\ref{conj:converse} to hold, and (3)---for a \emph{positive} outcome to
be a point of disagreement rather than one both predict---the clean
(state-functional) case of Section~\ref{sec:separation} to be realized,
which on current knowledge is itself conjectural; where any of these
fail (in particular, where a no-slicing interpreter invokes a birth
surface---available at no cost---or borrows the Lorentz-violating
completion) a \emph{positive}
detection no longer discriminates. The null direction still does,
however: by the asymmetry of Section~\ref{sec:separation} the beables
class \emph{entails} the ghost-sector signal (Conjecture~\ref{conj:converse})
while the no-slicing class merely permits it and retains a
covariance-preserving escape, so a null---at the template's sensitivity,
and for a sector the template's assumptions fit---would falsify the one
and not
the other---an experiment the two readings answer differently, which
already defeats the bookkeeping charge under the exposure conditions
(1)--(2). The clean state-functional case (3) is needed only for a
positive result to symmetrically \emph{select} the beables reading, not
to make the two readings empirically distinct. A bookkeeping dispute
with an empirical arbiter is not a bookkeeping dispute---and that
arbiter (a null at template sensitivity) exists wherever a
\emph{template-saturating} ghost sector does and
Conjecture~\ref{conj:converse} holds.

%=====================================================================
\section{The trilemma locates Everett}
\label{sec:trilemma}

The program's field-theory companions \cite{instantaneous,choosing}
assemble their results into a trilemma, which we restate in the form
argued there and will not strengthen here:

\begin{claim}[the ghost trilemma; tree level]
\label{claim:trilemma}
Coupled ghosts; \emph{empirical} Lorentz invariance (no
\emph{in-principle} observationally accessible preferred
foliation---``accessible'' meaning that some observable depends on the
foliation with nonzero effect, realized detectability being a separate,
flux-dependent question); definite
configurations (beables): any two, never all three, at tree level.
\end{claim}

The engine of the claim is the instantaneity result of
\cite{instantaneous}: at tree level, the Lorentz-invariant eternal
vacuum of a ghost coupled to ordinary matter is, in the stochastic
currency, an \emph{instantaneous state}---the boost-orbit volume
renders the decay multiplicity divergent, the minimal semigroup
degenerates, and a non-degenerate configuration-space process exists
only relative to a structure that individuates the decay channels: in
every known rescue, a frame or a birth surface. Every
configuration-preserving rescue in the literature (momentum cutoffs,
ghost condensation, a born vacuum) supplies such a structure, and
thereby---this is the other half---exposes it observationally through
the frame-locking of the decay products. Meanwhile the
definition-level rescues that keep Lorentz invariance, the
$\PT$/no-ghost quantizations, exit the sample space: the companion
\cite{choosing} proves a no-representation lemma---a real-eigenvalue
Jordan-block Hamiltonian admits no closed stochastic-process
representation under any similarity, a result that lifts to the full
Fock tower of a conformal mode---so those quantizations retain no
definite-configurations seat to sit in. Two honesty flags from the
source papers carry over unchanged in substance: the claim is \emph{tree level}
(a resummed suppression rendering the boost-orbit sum convergent
would relax the first horn), and it is \emph{argued against every
example the program could find, not proven exhaustive}.

Table~\ref{tab:horns} lays out the horns with the price each pays,
as computed inside the program.

\begin{table}[ht]
\centering
\small
\begin{tabular}{@{}p{0.24\linewidth}p{0.33\linewidth}p{0.34\linewidth}@{}}
\toprule
\textbf{Horn (what is dropped)} & \textbf{What is kept} & \textbf{The price, as priced in this program}\\
\midrule
Empirical Lorentz invariance & coupled ghosts $+$ definite
configurations & an observable preferred frame; on the
conformal-gravity instance the price forks on the gravitational
coupling $\alpha_g$, both branches costly \cite{choosing}\\[4pt]
Coupled ghosts & empirical Lorentz invariance $+$ definite
configurations & every ghost must be exorcised, and every exorcism is
a survival-conditioning with a computable cost \cite{exorcised}\\[4pt]
Definite configurations & coupled ghosts $+$ empirical Lorentz
invariance & no sample space: no closed stochastic process on
configurations exists to be had \cite{choosing}\\
\bottomrule
\end{tabular}
\caption{The three horns of the trilemma of
\cite{instantaneous,choosing}, with prices. Everett occupies the
third row.}
\label{tab:horns}
\end{table}

A word on each price. On the first horn, the frame is not merely
philosophical furniture: it is (conditionally, at tree level)
observable in the sector's decay
products (Section~\ref{sec:separation}), and on the flagship
conformal-gravity instance the companion paper's pricing forks on the
value of the dimensionless gravitational coupling---the face-value
rotation-curve branch is closed by solar-system constraints on the
published (and disputed) analyses, while the alternative branch
forces the regulating cutoff to macroscopic wavelengths
\cite{choosing}. On the second horn, keeping the world ghost-free is
itself a policy with a price schedule: the companion
\cite{exorcised} shows that each ghost-removal prescription on
record---a boundary condition, a contour, a norm---acts as a
survival-conditioning of the underlying process, with a closed-form
cost that in the flagship case collapses exactly at the point where
the exorcism is most needed. On the third horn the price is the
absence of a sample space: there is no stochastic process on definite
configurations at all, and hence no single actual history.

\paragraph{The placement.} Everett sits on the third horn---and, and
this is the point of this section, sits there \emph{comfortably}.
Dropping definite configurations is not something the trilemma does
\emph{to} Everett; it is what Everett walked in the door already
having done. The Everettian never wanted a sample space: probability
in the Emergent Multiverse enters through decision theory and
self-location, not through a stochastic process on configurations
\cite{Wallace2012}. And the reward for paying that price is exactly
the good the other horns forfeit: with no beables there is no beable
kinematics, no foliation for it to be anchored to, no frame for a
ghost sector to expose---manifest covariance survives coupled ghosts
unscathed. The trilemma therefore \emph{locates} Everett rather than
criticizing it: Everett is the pick-two solution
$\{$ghosts, Lorentz$\}$, and its famous virtue and its famous
vertigo are revealed as the two faces of a single purchase. The
trilemma does not embarrass any horn; it prices each.

An honesty paragraph, to keep the location claim from quietly
inflating into an argument. The instantaneity theorem of
\cite{instantaneous} is \emph{beables-relative}: it says a
configuration-space process for a coupled ghost vacuum degenerates
unless individuating structure is supplied. Everett has no
configuration-space process, so the theorem does not bite Everett at
all---which is just the location claim restated from the other side.
It is worth being explicit that the trilemma locates Everett by
\emph{negation}: nothing in the trilemma's mechanism, which operates on
configuration-space processes, reaches Everett at all, and it is
Everett's \emph{own} commitment to having no beables that seats it on
the drop-configurations horn. The placement is done by Everett's
premise, not by any argument of ours. A further dependency we state
once and plainly: claims (ii) and (iii) inherit their content from the
program's companion papers \cite{instantaneous,choosing,exorcised}---%
the trilemma, the instantaneity theorem, the exposure template, and the
exorcism costs are established there, not here---so claim~(ii)'s
soundness is exactly the soundness of those companions, while
claim~(iii) rests on them \emph{and}, over and above, on
Conjecture~\ref{conj:converse} (with the comoving-frame and clean-case
conditions of Section~\ref{sec:separation}), this paper's own and
admittedly conjectural; this paper assembles and prices their results,
it does not reprove them. Nothing in this section
is evidence against the Everett interpretation, and we do not intend a
cumulative-case reading in which a trilemma horn is a demerit. Horns
are addresses, not verdicts.

%=====================================================================
\section{The separation: a ghost sector makes it empirical}
\label{sec:separation}

This is the paper's central section and its one novel claim. The
claim is conditional and one of its two supports is a conjecture; we
build it in three steps and flag the status of each.

\paragraph{Step 1 (companion results: theorem at the core, argument
at field strength).} A beables-first field theory is
foliation-anchored by construction: defining a configuration space
for field beables requires a choice of spacelike slicing, exactly as
in Bohmian field theory \cite{DGNSZ}. The companion paper
\cite{instantaneous} sharpens this ancient observation in two ways
that matter here. First, the foliation's \emph{determinacy} is a
well-defined (quantifiable) functional of the quantum state: the paper
exhibits a measurable property of the state quantifying
how sharply the state selects a frame, and shows---by exact algebra at
the vacuum, and a measured determinacy curve approaching it---that it
\emph{vanishes continuously at the vacuum}. A perfect-fluid state with
$\rho+p\neq 0$ determines a frame; the vacuum determines none; the
boundary $\rho+p=0$ is the phantom divide. Second---the shield---in
ordinary (ghost-free) sectors in quantum equilibrium the anchoring
foliation is empirically hidden: no Born-rule observable depends on
it. Within the indivisible framework this hiddenness is proven as a
no-signaling theorem for the finite-configuration core, and
\emph{conjectured}, with a defined proof target, at full field-theory
strength \cite{instantaneous}.

It is this shield that makes our position an extension of Barandes'
rather than a heresy within his framework. Barandes' published claims
that the formulation is causally local \cite{BarandesCausal,BHK2025}
and has no fundamentally preferred foliation \cite{BarandesPoP} concern
the causal
and signaling structure of quantum theory as instantiated in ordinary
matter---and there we agree, on the strength of the same shield: on
every sector so far observed, the anchoring foliation is
gauge, exactly as unobservable as his theorems require. What follows
concerns a sector nobody has observed. We should flag, though, that
this reconciliation is our reading and not Barandes'. Restricting ``no
fundamentally preferred foliation'' to ordinary matter, and treating a
ghost sector as the regime where the foliation is promoted to
observable, is a substantive interpretive move on which Barandes may
well disagree: his stated position is that the formulation has no
fundamentally preferred foliation, without a matter-content proviso. We
claim compatibility with the letter of his theorems, which concern
observed sectors, not his endorsement of the extension.

\paragraph{Step 2 (this paper's step, and it is conjectural).} A
coupled ghost sector would break the hiding. The mechanism is the
first horn's exposure result run in reverse: a coupled ghost's decay
products are isotropic in the individuating frame and in no other,
so a detector moving with velocity $\beta$ relative to that frame
sees a Compton--Getting kinematic dipole aligned with (on the
natural identification of the individuating frame with the comoving
frame) the CMB dipole, together with a direction-modulated spectral
edge $\Lambda(\theta)\simeq\Lambda(1+\beta\cos\theta)$ at the
sector's cutoff. The companion paper computed the template
quantitatively for the phantom instance \cite{instantaneous}:
count-rate dipole amplitudes of a few$\,\times10^{-3}$ at
$\beta=1.23\times10^{-3}$, roughly $4.0$--$5.5\times10^{6}$ signal
photons for a $5\sigma$ dipole detection with the apex fixed by the
CMB dipole, and an edge modulation at the $10^{-3}$ level recoverable
through realistic instrument resolution. The discriminant against
mundane comoving backgrounds, which share the Compton--Getting
dipole, is the edge modulation together with the absence of source
correlation. One load-bearing step here is an assumption, not a
derivation, and we flag it as such: the identification of the
individuating frame with the comoving (CMB) frame. The instantaneity
result of \cite{instantaneous} fixes only that \emph{some} frame
individuates the decay channels; that this is the frame in which the
cosmic matter distribution is at rest is an assumption we make---%
plausible for a cosmologically prepared sector---but do not derive.
It is worth being explicit about which rescue underlies this
quantitative template: it is computed for a \emph{frame-cutoff}
instance---a momentum cutoff imposed in the individuating frame, which
we identify above with the comoving (CMB) frame---rather than for a
purely state-functional foliation. The Compton--Getting dipole is
generic to any individuating frame, but the spectral edge, and the
saturation bound $\Lambda\lesssim19\,$MeV that makes the photon counts
feasible, both presuppose that cutoff. By the analysis of Step~3 below,
the template therefore already sits in the \emph{confounded} branch;
and because a momentum cutoff is a law-level Lorentz-violating
completion, it sits in the branch a no-slicing interpreter may adopt
only at the price of manifest covariance, not in the birth-surface
branch available for free. It should not be read as a signature of the
clean case. If
the individuating frame were some other frame, the dipole apex would
move with it and the CMB-locked template would need re-pointing; a
frame-locked signal would still exist, but its specific direction would
not be predicted. We emphasize that these numbers live \emph{inside}
the conditional built in this section: they quantify the template that
becomes operative only if a coupled ghost sector exists and the
conjecture below holds.

We state the logical status of this step with maximal bluntness,
because it is the paper's load-bearing conditional. The observability
claim is the \emph{converse} of the shield: the shield says ordinary
sectors hide the foliation; the converse says a ghost sector exposes
it. In \cite{instantaneous} the shield itself is a theorem only for
the finite-configuration core, and the ghost-induced observability is
argued at tree level, not proven. If the full-strength shield were to
fail, the separation below would be confounded (foliation effects
might leak elsewhere); if the tree-level exposure argument were to
fail under resummation, the signal might be absent even on the
foliation-anchored reading. We therefore record the step as a
conjecture, whose resolution is the natural next proof target for the
program:

\begin{conjecture}[shield converse; conditional]
\label{conj:converse}
In a beables-first indivisible field theory whose matter content
includes a ghost sector coupled to ordinary matter, some Born-rule
observable of the ghost sector's decay products depends on the
anchoring foliation with nonzero effect (generically---away from the
phantom divide $\rho+p=0$, where the foliation-determinacy functional
vanishes)---the kinematic dipole and
spectral edge of \cite{instantaneous} being the template---while all
observables of ghost-free sectors remain foliation-independent.
\end{conjecture}

Two conjuncts are bundled here for economy, but only the first is
load-bearing for the falsification below. The null test uses the
\emph{exposure} conjunct alone---beables $\Rightarrow$ a ghost-sector
signal, so a null $\Rightarrow$ not beables, as the template realizes
it (Conjecture~\ref{conj:converse} delivers a \emph{nonzero} effect;
reaching a null's full force needs the template's saturating-cutoff and
comoving-frame auxiliaries, Step~2 above); the second conjunct (the
\emph{shield} for ghost-free sectors, already a theorem for the finite
core) governs only whether foliation effects might leak into ordinary
sectors and confound the test, not whether the null falsifies. Keying
the falsification to exposure alone keeps it from being hostage to the
field-strength shield.

\paragraph{Step 3 (the separation, and where its positive direction is
confounded).} A no-preferred-slicing ontology predicts no such signal
\emph{when the ghost Lagrangian is held fixed and Lorentz-invariant and
no individuating completion or birth-surface initial condition is
invoked}, because then it has nothing for the signal to reveal. On the Everettian
reading of the same Lorentz-invariant field content there is no beable
kinematics, hence no anchoring foliation, hence no frame for the decay
products to lock to; the natural Everettian completions of such a ghost
sector are precisely the definition-level, covariance-preserving
quantizations that the trilemma's third horn comprises---and those have
no frame-locked decay signal to give.

The positive direction, however, carries a confound we must state
plainly, because the trilemma's own exposure mechanism supplies it. In
Section~\ref{sec:trilemma} the individuating structure that rescues a
coupled ghost vacuum is furnished, in every rescue the program could
find, by one of two kinds of structure, and it is worth keeping them
apart. The first is a physically \emph{Lorentz-violating} completion at
the level of the law---a momentum cutoff or ghost condensation. The
second is a \emph{birth-surface} (initial-condition) structure such as
a born vacuum, which is not a law-level Lorentz violation at all but a
cosmological initial datum; though sometimes grouped with the
Lorentz-violating rescues, it is more accurately classed as an
initial condition. Both routes are available to a no-slicing
(Everettian) interpreter. The Lorentz-violating completion can simply
be adopted; and the birth surface is available \emph{for free} to any
Everettian who accepts big-bang cosmology, so if a bare cosmological
initial condition suffices to individuate the decay channels, the
confound is \emph{generic}---essentially every realistic Everettian
can then reproduce the same signal at no cost, and the positive
direction is weaker still. Either way such an interpreter can adopt the identical structure
and predict the \emph{same} frame-locked dipole and edge: in that case
the signal tracks the completion or the initial condition, not the
ontology. A positive detection therefore does \emph{not} entail a
foliation-anchored ontology. The
separation is clean only in the restricted case where the individuating
frame arises \emph{solely} as a state-functional foliation---the
covariant-foliation-from-the-wave-function picture of \cite{DGNSZ},
extracted from the state with the Lagrangian held fixed and
Lorentz-invariant---rather than from a fundamentally Lorentz-violating
completion or a birth surface. Wherever such a completion or initial
condition is in play, a positive signal is confounded between the two
classes. This clean case should therefore be named as a \emph{third,
separate condition} on the separation, alongside the existence of a
coupled ghost sector and Conjecture~\ref{conj:converse}: that a
state-functional foliation \emph{suffices} to individuate the decay
channels. And we must concede that this condition has never been met on
paper. No purely state-functional rescue of a ghost vacuum has ever
been exhibited: by this program's own survey
(Section~\ref{sec:trilemma}), every known rescue supplies a frame or a
birth surface---a momentum cutoff, ghost condensation, a born
vacuum---and the \cite{DGNSZ} construction is a picture for ordinary
Bohmian matter, not a demonstrated rescue of the \emph{ghost} vacuum.
Conjecture~\ref{conj:converse}, whose antecedent presupposes only a
beables-first ghost field theory, does not cover this gap. On current
knowledge, therefore, every concrete realization of the positive signal
lies in the confounded branch, and the clean case is conjectural
\emph{over and above} Conjecture~\ref{conj:converse}. Moreover the
birth-surface route is not a gap that further modeling might simply
close: an FRW birth surface is generic to any big-bang cosmology, so
\emph{if} a bare cosmological initial condition suffices to individuate
the decay channels, the confound on the positive direction is not
merely a present lacuna but arguably \emph{structural}---unobservable in
principle on a realistic cosmological background, not just
so-far-unrealized. This strengthens,
rather than weakens, the null as the discriminator: if every
\emph{viable} beables completion must carry an individuating structure,
then every viable beables completion predicts a signal, so the absence
of one \emph{at the template's sensitivity} presses hard against every
template-saturating realization of the class. The
discriminating power is thus asymmetric, and the asymmetry is what
keeps claim~(iii) non-vacuous. Given Conjecture~\ref{conj:converse} and
the program's survey, the beables class \emph{entails} a nonzero signal,
so for an established sector whose computed template puts that signal at
the null's sensitivity a null result \emph{falsifies} that class \emph{as
the template realizes it}---a nonzero but sub-threshold effect (a sector
the template does not cover, \emph{or} a covered sector realized off the
comoving frame or below the template's signal fraction) would survive it. The Everettian, by contrast,
merely \emph{permits} the signal: the quantization is a free choice---%
covariance-preserving with a Lorentz-invariant Lagrangian and no
completion, or signal-matching by borrowing a completion or birth
surface---so Everett entails nothing and remains compatible with
\emph{both} outcomes. A positive detection therefore refutes nothing on
the no-slicing side, whereas a null---at the template sensitivity, for an
established sector---falsifies the beables class (under
Conjecture~\ref{conj:converse}); only the beables direction carries an
entailment for an experiment to break.

Even a confounded positive, moreover, is not always without force for
the no-slicing side---but here we must scope the point carefully,
because the two confounding routes are not on a par. A no-slicing
interpreter who adopts a fundamentally \emph{Lorentz-violating}
completion at the level of the law in order to match the signal thereby
forfeits \emph{manifest covariance}---which
Sections~\ref{sec:intro} and~\ref{sec:trilemma} identify as exactly the
good that dropping configurations was supposed to buy. In that branch a
confounded positive fails to select the ontology, yet it still exacts a
real and nontrivial price from the horn whose entire purchase was
covariance. The birth-surface branch is different, and we concede it:
a Lorentz-invariant law with a non-invariant cosmological initial state
(an FRW birth surface) forfeits no manifest covariance at the level of
the law, so there the individuating structure is available \emph{for
free}---as noted above---and a confounded positive costs the Everettian
nothing. This law-level covariance does not, however, put a born-vacuum
beables theory in the trilemma's forbidden ``all three'' corner: the
trilemma's second element is \emph{empirical} Lorentz invariance---no
observationally accessible preferred foliation
(Claim~\ref{claim:trilemma})---and a born vacuum (conditionally, at tree
level) exposes the frame observationally, so a beables theory with a
Lorentz-invariant law but a frame-exposing birth surface still drops the
empirical-Lorentz-invariance horn. The law-level covariance it retains
is not the currency the trilemma trades in. The one branch so far
realized as a concrete signal template is, however, the
\emph{frame-cutoff} branch, which \emph{is} a law-level
Lorentz-violating completion; so the covariance price is one the
exhibited realization actually charges. It is the birth-surface
branch---the branch that would exact no such price, and whose
availability for free is what makes the confound generic---that has
not so far been realized as a concrete template. Putting the
steps together, conditionally, and with the confound flagged:

\begin{quote}
\emph{If} nature contains a coupled ghost sector, \emph{and if}
Conjecture~\ref{conj:converse} holds, \emph{then}: a frame-locked
dipole-plus-edge signal in that sector's decay products is
\emph{consistent with}---but does not select---the foliation-anchored
single-world class, being \emph{confounded} by any fundamentally
Lorentz-violating completion or
birth-surface initial condition of the sector, either of which
reproduces the same signal on a no-slicing ontology---the birth surface
at no cost, the Lorentz-violating completion at the price of manifest
covariance; the
signal indicates foliation-anchoring cleanly only under the further,
so-far-unrealized condition that the individuating frame arises purely
as a state-functional foliation with the Lagrangian held
Lorentz-invariant. A null result, with the ghost sector independently
established, is the cleaner discriminator, precisely because the
entailment runs one way: the beables class \emph{entails} a signal (of
the computed template amplitude, for an established sector, under the
assumed comoving-frame identification---a realization off that frame,
below the template's signal fraction, or individuated by a completion
other than the template's frame cutoff, survives the pointed null),
so its absence is evidence against definite configurations and presses
toward the no-slicing class, whereas a no-slicing ontology merely
\emph{permits} the signal, so a positive leaves it untouched.
\end{quote}

Neither outcome selects an interpretation. The positive outcome does
not distinguish Bohmian from stochastic beables
(Section~\ref{sec:limits}); the null outcome does not distinguish
Everett from any other foliation-free reading. What the experiment
discriminates is a \emph{class boundary}---slicing-anchored versus
not---which is, we believe, precisely the boundary the interpretation
debate has always treated as empirically inaccessible. And the Born
rule is intact on both sides throughout: the two classes differ not
about any probability rule but about what the world contains for the
rule to apply to, and the disagreement becomes observable only
because a ghost sector's kinematics is (conditionally) sensitive to
the difference.

\paragraph{The ancestry, and the contrast that defines the claim.}
The idea that a hidden-variable ontology might escape the
``Everett in denial'' charge by empirical divergence is Valentini's:
in quantum non-equilibrium, pilot-wave theory diverges observably
from quantum mechanics, so the corpuscles are not idle wheels
\cite{Valentini0811,Valentini1991a,Valentini2002}. But Valentini's
escape requires the Born rule to \emph{fail}---non-equilibrium is by
definition a violation of Born statistics, and his signal-locality
theorem ties the divergence to that violation. Ours is the
ghost-sector generalization of the same escape with the requirement
inverted: the discriminator is a sector, not a rule failure. One other
equilibrium route to foliation observability has been proposed---%
Bohmian arrival-time statistics in EPRB configurations
\cite{AmblardDrezet}, a laboratory-scale mechanism resting on the
contested measurability of arrival times, and disjoint from the
sector-based discriminator here. Beyond it we know of no prior proposal
that separates interpretation classes observationally while holding the
Born rule fixed. The nearest
relatives sharpen the contrast. D\"urr, Goldstein, Norsen, Struyve,
and Zangh\`i showed that a Bohmian theory can extract a covariant
foliation from the wave function itself, adding that such a foliation
is undetectable in quantum equilibrium \cite{DGNSZ}---our scenario is
precisely the loophole in that equilibrium hiding: it is a fact about
ordinary sectors, and a ghost sector
(conditionally) breaks it while equilibrium holds. Sahin's recent
program of quantum-field-theoretic phenomenology for
$\psi$-epistemic models \cite{Sahin2026} needs deviations from
quantum-mechanical predictions to generate signals; we need none.
And the fragility results of \cite{MMO2025} are companions in spirit:
empirical equivalence between interpretations is a regime-bound
theorem, not a decree, and our claim is that the regime has a
physically meaningful boundary.

\paragraph{The conditional, stated plainly.} The trigger may never
fire. If nature contains no coupled ghost sector, the separation
never becomes operative and the equivalence between the classes
stands unbroken---this paper's claim is a conditional falsifiability
statement, the same logical DNA as the rest of the corpus. What would
count: on the input side, an established phantom sector (a
$w<-1$ equation of state confirmed jointly with structure-growth data
discriminating a genuine field from a gravitational mimic) is the
scenario in which the template of \cite{instantaneous} applies as
computed; current cosmological data contain disputed hints and
nothing more, and we take no position on them here, deferring to the
companion's discussion. On the output side, the null is well-defined:
ghost sector established, exposure and systematics sufficient for
the $10^{-3}$ dipole and edge, and no frame-locked signal. We
emphasize once more that the detection numbers quoted above sit on a
knife-edge of signal fraction detailed in \cite{instantaneous}; the
separation is an experiment nature may offer, not one we can
schedule. Absent the sector, the statement is vacuously safe and
empirically idle---and we say so.

%=====================================================================
\section{What does not follow}
\label{sec:limits}

Four disclaimers, each of which we regard as part of the paper's
content rather than throat-clearing.

\emph{(i) The instantaneity theorem is not an argument against
Everett.} It is beables-relative (Section~\ref{sec:trilemma}): it
constrains configuration-space processes, and Everett has none. We
have used it only to price the horn Everett declines, never to argue
against the horn Everett occupies; any reading of this corpus as a
cumulative case against the Everett interpretation is a misreading.

\emph{(ii) Division events are not branch events---so far as anyone
has shown.} There is a tempting translation: division events, the
isolated times at which an indivisible process factorizes, resemble
the decoherence times at which Everettian branch structure
stabilizes; and the terminal-indivisibility class of
\cite{indivisible}---no division events beyond a finite time---would
translate to ``no stable branch structure available in ghost
regimes.'' We flag this correspondence as \emph{suggestive
translation only}. The program's theorems concern
configuration-level divisors; Everettian branching is defined through
decoherence functionals on coarse-grained histories; no theorem
connecting the two exists, in this corpus or elsewhere to our
knowledge. It is a conjecture-shaped hole, not a claim.

\emph{(iii) Nothing here selects among single-world readings.} The
separation of Section~\ref{sec:separation} discriminates
slicing-anchored from non-slicing ontologies. A positive signal would
not distinguish Bohmian beables from stochastic beables, nor either
from any other foliation-anchored ontology; a null would not
distinguish Everett from foliation-free alternatives. The experiment,
if nature grants it, adjudicates a class boundary and nothing finer.

\emph{(iv) The separation claim is conjectural several times over.} It
conditions on the existence of a coupled ghost sector (an empirical
question, open), and on Conjecture~\ref{conj:converse} (a theoretical
question, open: the shield is proven only for the finite core, and
the exposure argument is tree-level); and its \emph{clean} (class-%
selecting) direction conditions further on a state-functional foliation
sufficing to individuate the decay channels---a condition that, on
current knowledge, no exhibited rescue meets, so the confounded branch
is the only one so far realized. We have tried to write every
sentence of Section~\ref{sec:separation} so that it survives the
worst case on all three conditions; where a sentence seems to outrun
them, the conditions govern.

%=====================================================================
\section{Discussion}
\label{sec:discussion}

Interpretations of quantum mechanics that differ in a register
physics can reach are rare; that has always been the debate's
melancholy. The argument of this paper is that the register exists
and has a name: ghosts. A coupled ghost sector is the one place the program has found in
field theory where, at tree level, the existence of a sample space
stops being free---where definite configurations demand an
individuating structure and (conditionally) advertise it, and where
declining to have configurations converts the signal from an entailment
into a mere permission. If the reconstructions of
Section~\ref{sec:landscape} restore the underdetermination that the
no-alternative argument had claimed to close, the trilemma and the
separation say where the underdetermination could end: not in the
seminar room, but in the direction-dependence of a diffuse
background.

We do not expect the question settled soon, and the paper's
architecture is built for that: claim (i) stands as a piece of
underdetermination independent of any ghost physics---with two
qualifications this paper has been at pains to make, that it is
conditional on Barandes' (cited, not here established) exit from the
PBR theorem, and that its \emph{defense} against the strongest
Everettian reply, the bookkeeping objection, is not fully
ghost-independent, since the decisive answer to that reply is the
ghost-dependent claim (iii); unconditionally, claim (i) meets the
objection only with the defensive parity and disanalogy replies.
Claim (ii) stands at tree level whatever the observations do; claim
(iii) waits on a sector, a conjecture, and---for its positive
direction---an unexhibited clean case. What the trilemma does
unconditionally is raise the stakes of the ghost question itself. Whether nature contains coupled ghosts has been, to
date, a question about instability and regularization; on the
argument here it is also, conditionally, the question of whether the
interpretation debate has an empirical future.

We end with the two vertigos, because the corpus should take no side
the data cannot take. The Everettian asks us to accept worlds beyond
counting, splitting without cease, and offers in exchange a formalism
taken perfectly literally. The stochastic reading asks us to accept
memory without storage---a world whose single history obeys a law
that remembers what no beable records---and offers in exchange a
single world. Both offers are coherent; both prices are real; and
the argument of this paper is that, if nature is unkind enough to
contain a ghost sector, she may also be kind enough to tell us which
price she has already paid.

\paragraph{Author contributions and use of AI.} This manuscript was
researched and written by the large language model Claude Fable~5
(Anthropic), working under the direction of Jonathan Torkelson, who
conceived the program's originating question, directed the sessions,
and reviewed the outputs. The author is a professional engineer, not
a physicist, and makes no claim to independent verification of the
technical content; the manuscript is circulated precisely to obtain
such verification from professional physicists. All errors are owned
by the process that produced them.

\paragraph{Code availability.} This is a synthesis paper and proves
no new theorem and runs no new numerics. Every numerical claim quoted
here---the divisor computations and terminal-indivisibility
constants, the dipole and spectral-edge template and its photon
counts, the foliation-determinacy functional---is inherited from the
companion papers and reproducible from their public scripts; entry
points and the program-wide index are collected at the hub repository
\url{https://github.com/embeddetech/ghosts-program}.

%=====================================================================
\begin{thebibliography}{99}\small

\bibitem{indivisible} J.~Torkelson,
\emph{Indivisible stochastic processes for non-Hermitian and
higher-derivative quantum systems: detailed balance, broken $\PT$
symmetry, and the explosion of ghosts} (2026),
\url{https://github.com/embeddetech/p01-indivisible-ghosts}.

\bibitem{instantaneous} J.~Torkelson,
\emph{Ghost vacuum decay as a stochastic process: instantaneous
states, the boost-orbit measure, and the foliation that ghosts make
visible} (2026),
\url{https://github.com/embeddetech/p03-instantaneous-ghosts}.

\bibitem{choosing} J.~Torkelson,
\emph{Conformal gravity must choose: exceptional-point kinematics,
vacuum decay, and the sample space of Weyl gravity} (2026),
\url{https://github.com/embeddetech/p05-choosing-ghosts}.

\bibitem{exorcised} J.~Torkelson,
\emph{Every exorcism has a price: ghost removal as
survival-conditioning, from Maldacena's boundary condition to the
pole-sheet divisor} (2026),
\url{https://github.com/embeddetech/p08-exorcised-ghosts}.

\bibitem{BarandesPoP} J.~A.~Barandes,
\emph{The Stochastic-Quantum Correspondence},
Philosophy of Physics \textbf{3}(1), 8 (2025),
doi:10.31389/pop.186, arXiv:2302.10778.

\bibitem{BarandesTheorem} J.~A.~Barandes,
\emph{The Stochastic-Quantum Theorem}, arXiv:2309.03085 (2023).

\bibitem{BarandesCausal} J.~A.~Barandes,
\emph{New Prospects for a Causally Local Formulation of Quantum
Theory}, arXiv:2402.16935 (2024).

\bibitem{BHK2025} J.~A.~Barandes, M.~Hasan, and D.~Kagan,
\emph{The CHSH Game, Tsirelson's Bound, and Causal Locality},
arXiv:2512.18105 (2025).

\bibitem{BarandesHMM} J.~A.~Barandes,
\emph{Pilot-Wave Theories as Hidden Markov Models},
arXiv:2602.10569 (2026).

\bibitem{HS2010} N.~Harrigan and R.~W.~Spekkens,
\emph{Einstein, incompleteness, and the epistemic view of quantum
states}, Found.\ Phys.\ \textbf{40}, 125 (2010), arXiv:0706.2661.

\bibitem{PBR} M.~F.~Pusey, J.~Barrett, and T.~Rudolph,
\emph{On the reality of the quantum state},
Nature Phys.\ \textbf{8}, 475 (2012), arXiv:1111.3328.

\bibitem{PBRnoPI} S.~Gao,
\emph{From joint to single-system psi-onticity without preparation
independence}, Found.\ Phys.\ \textbf{56}, 7 (2026),
arXiv:2601.17662.

\bibitem{Deutsch1996} D.~Deutsch,
\emph{Comment on Lockwood},
Brit.\ J.\ Phil.\ Sci.\ \textbf{47}, 222--228 (1996);
the quoted sentence is at p.~225.

\bibitem{BrownWallace} H.~R.~Brown and D.~Wallace,
\emph{Solving the measurement problem: de Broglie--Bohm loses out to
Everett}, Found.\ Phys.\ \textbf{35}, 517 (2005),
arXiv:quant-ph/0403094.

\bibitem{Valentini0811} A.~Valentini,
\emph{De Broglie--Bohm pilot-wave theory: many worlds in denial?},
in \emph{Many Worlds? Everett, Quantum Theory, and Reality},
Oxford University Press (2010), arXiv:0811.0810.

\bibitem{Brown0901} H.~R.~Brown,
\emph{Reply to Valentini, ``De Broglie--Bohm pilot-wave theory: many
worlds in denial?''}, in \emph{Many Worlds? Everett, Quantum Theory,
and Reality}, Oxford University Press (2010), pp.~510--518,
arXiv:0901.1278.

\bibitem{Wallace2012} D.~Wallace,
\emph{The Emergent Multiverse: Quantum Theory according to the
Everett Interpretation}, Oxford University Press (2012).

\bibitem{Valentini1991a} A.~Valentini,
\emph{Signal-locality, uncertainty, and the subquantum H-theorem.~I},
Phys.\ Lett.\ A \textbf{156}, 5 (1991).

\bibitem{Valentini2002} A.~Valentini,
\emph{Subquantum information and computation},
Pramana \textbf{59}, 269 (2002), arXiv:quant-ph/0203049.

\bibitem{DGNSZ} D.~D\"urr, S.~Goldstein, T.~Norsen, W.~Struyve, and
N.~Zangh\`i,
\emph{Can Bohmian mechanics be made relativistic?},
Proc.\ R.\ Soc.\ A \textbf{470}, 20130699 (2014), arXiv:1307.1714.

\bibitem{AmblardDrezet}
A.~Amblard and A.~Drezet,
\emph{Arrival-time distributions as a probe of the preferred
foliation in relativistic Bohmian mechanics},
arXiv:2604.17507.

\bibitem{MMO2025} J.~Manero, R.~Muci\~no, and E.~Okon,
\emph{A disputable assumption behind the empirical equivalence
between pilot-wave theory and standard quantum mechanics},
arXiv:2508.06667 (2025).

\bibitem{Sahin2026} I.~Sahin,
\emph{Quantum field-theoretic predictions of $\Psi$-epistemic models
of quantum mechanics}, arXiv:2605.12546 (2026).

\bibitem{BenderMannheim2008} C.~M.~Bender and P.~D.~Mannheim,
\emph{No-ghost theorem for the fourth-order derivative
Pais--Uhlenbeck oscillator model},
Phys.\ Rev.\ Lett.\ \textbf{100}, 110402 (2008), arXiv:0706.0207.

\bibitem{Mannheim2020} P.~D.~Mannheim,
\emph{Ghost problems from Pauli--Villars to fourth-order quantum
gravity and their resolution},
Int.\ J.\ Mod.\ Phys.\ D \textbf{29}, 2043009 (2020).

\end{thebibliography}

\end{document}
